% LỜI TỰA --- Quyển IV
\chapter*{Lời Tựa}
\addcontentsline{toc}{chapter}{Lời Tựa}
\markboth{Lời Tựa}{Lời Tựa}

\begin{truyen}{Lời Tựa}{Tales of Life Philosophy}
\chuhoa{G}{iáo viên bắt đầu bộ sách này với một câu hỏi đơn giản: làm thế nào để mỗi buổi học có thể để lại thứ gì đó đáng nhớ hơn là bài thi?}
\textit{(A teacher began this series with a simple question: how can each class leave behind something more memorable than a test?)}

Triết lý nhân sinh không chỉ là của các nhà triết học~--- nó ở ngay trong những chuyện đời thường.
\textit{(Life philosophy is not just for philosophers~--- it lives right inside everyday stories.)}

Quyển bốn đưa người đọc đến với những tình huống quen thuộc từ cổ chí kim: một người thầy thay đổi học trò, một lời khuyên đúng lúc thay đổi số phận, một bài học từ thất bại mang lại thành công. Triết lý nhân sinh không phải để tranh luận~--- nó để sống. Mỗi câu chuyện trong quyển này là một minh chứng cho điều đó.
\textit{(Volume four takes the reader through familiar situations from ancient to modern times: a teacher changing a student, timely advice changing a fate, a lesson from failure leading to success. Life philosophy is not for debate~--- it is for living. Each story in this volume is evidence of that.)}

Bộ sách <<Truyện Tích Cảm Hứng Song Ngữ>> gồm 24 quyển, mỗi quyển một chủ đề riêng. Quyển IV này mang chủ đề: \textbf{triết lý nhân sinh từ cổ chí kim}.
\textit{(The series ``Inspiring Bilingual Tales'' contains 24 volumes, each with its own theme. Volume IV focuses on: \textbf{triết lý nhân sinh từ cổ chí kim}.})
\end{truyen}

\begin{baihoc}
Người trí không chỉ học từ kinh nghiệm bản thân~--- mà còn học từ những câu chuyện đã xảy ra trước mình.
\textit{(The wise learn not only from their own experience~--- but also from stories that happened before them.)}
\end{baihoc}

\begin{tcolorbox}[
  colback=truyenblue!6, colframe=truyenblue,
  coltitle=white, fonttitle=\bfseries,
  title={Easy English},
  arc=4pt, boxrule=1pt, breakable, enhanced
]
\textbf{Short English to remember:}

Ancient wisdom speaks to modern lives.

Every story is a lesson waiting to be lived.
\end{tcolorbox}

\begin{tcolorbox}[
  colback=truyengold!10, colframe=truyengold!70!black,
  coltitle=black, fonttitle=\bfseries,
  title={T\textit{ự} H\textit{ọ}c Nhanh},
  arc=4pt, boxrule=1pt, breakable, enhanced
]
1. Triết lý nhân sinh nào trong quyển này em thấy đúng nhất với cuộc sống hiện tại? \textit{(Which life philosophy in this volume rings truest for your life right now?)}

2. Câu chuyện nào thay đổi cách em nhìn một vấn đề? \textit{(Which story changed the way you see a problem?)}

3. Em học được gì từ những sai lầm và vấp ngã được kể trong quyển này? \textit{(What did you learn from the failures and stumbles described in this volume?)}
\end{tcolorbox}

\begin{center}
\textit{\color{truyengold}``Từ những câu chuyện xưa, ta tìm thấy chính mình.''}
\end{center}

\begin{flushright}
\textbf{Tôi Nguyễn Văn Sang}\\
\textit{GV THPT Nguyễn Hữu Cảnh - TP HCM}
\end{flushright}

\ngancach
